% !PRINTED 214
\weberpage{214}
%<<CONT
which, in view of $\Theta^2 = \psi^2 = \chi^3 = 1$, follow from the three
relations
\[
\psi\Theta = \Theta\psi, \quad \chi\Theta = \Theta\psi\chi, \quad \chi\psi = \Theta\chi
\]

From (5) it further follows that the four-group
\weq{6}{1,\ \Theta,\ \psi,\ \Theta\psi}
is a \gesperrt{normal divisor} of the tetrahedral group.

The substitutions of the tetrahedral group can, in each of the
two forms
\[
\chi^\lambda \psi^\mu \Theta^\nu, \quad \Theta^\nu \psi^\mu \chi^\lambda
\]
\[
\lambda = 0, 1, 2, \quad \mu = 0, 1, \quad \nu = 0, 1
\]
be represented.

To obtain the two fundamental forms of the fourth degree $\Phi_1, \Phi_2$
belonging to the three-fold poles, one can either compute
from (5) the poles still missing, or one can proceed as
follows:

Since the fundamental forms $\Phi_1, \Phi_2$ must admit the substitutions $\Theta$ and $\psi$,
and cannot be divisible by $x_1$ and $x_2$,
they can only be of the following form:
\[
\begin{aligned}
\Phi_1 &= x_1^4 + m_1 x_1^2 x_2^2 + x_2^4 \\
\Phi_2 &= x_1^4 + m_2 x_1^2 x_2^2 + x_2^4,
\end{aligned}
\]
in which $m_1, m_2$ are constants still to be determined. But if one forms
the \gesperrt{Hessian} covariant of $\Phi_1$, one obtains by
\S.\,40, 4. an invariant of the group:
\[
\frac{1}{12}\left[\frac{\partial^2 \Phi_1}{\partial x_1^2}\,\frac{\partial^2 \Phi_1}{\partial x_2^2} - \left(\frac{\partial^2 \Phi_1}{\partial x_1\,\partial x_2}\right)^2\right] = 2m_1(x_1^4+x_2^4) + (12-m_1^2)\,x_1^2 x_2^2.
\]

This form, since no other fundamental forms of the fourth
degree exist, can agree only with $\Phi_1$ or with $\Phi_2$
(apart from a constant factor). Hence we must have
\[
\frac{12-m_1^2}{2m_1} = m_1 \ \text{or}\ m_2
\]
The first assumption, however, gives $m_1 = \pm 2$. Then $\Phi_1$
would be a square, which cannot be. There remains, then, only:
\[
m_1^2 + 2m_1 m_2 = 12,
\]
and likewise from $\Phi_2$:
\[
m_2^2 + 2m_1 m_2 = 12,
\]
hence
\[
m_1 = -m_2 = \pm 2\sqrt{-3},
\]
%CONT>>
